\section[Exact algebra for the 3-by-2 separation]
{Exact algebra for the \(3\times2\) separation}
\label{app:separation-algebra}

We record the identities that make the projective bound independent of
numerical optimization. For a pair involving auxiliary label \(2\), let
\[
A=\frac35D \Pi(\bm n)D,\qquad
B=\frac45D \Pi(\bm m)D,
\]
where \(x=n_z\), \(y=m_z\), \(r=\bm n\cdot\bm m\), and
\[
D=\operatorname{diag}\!\left(
\sqrt{\frac{1+\eta}{2}},\sqrt{\frac{1-\eta}{2}}\right).
\]
Direct expansion, using \(\|\bm n\|=\|\bm m\|=1\), gives
\[
\det(A-B)=-\frac3{50}(1-\eta^2)(1-r)\le0.
\]
Thus \(A-B\) has eigenvalues of opposite signs, including a possible zero
boundary eigenvalue, and the optimized complementary-projector score is
\[
F_\eta=\frac{\Tr A+\Tr B+\|A-B\|_1}{2}.
\]
Its discriminant yields
\[
20F_\eta=
7+\eta(3x+4y)+
\sqrt{[-1+\eta(3x-4y)]^2+24(1-\eta^2)(1-r)}.
\]

At \(\eta=0\),
\[
F_0(r)=\frac7{20}+\frac1{20}\sqrt{25-24r}.
\]
Set
\[
q=\sqrt{25-24r}\ge1,\qquad H=8-3x-4y\ge1.
\]
To prove \(F_\eta\le F_0(r)+2\eta/5\), it is enough to prove
\[
\sqrt{[-1+\eta(3x-4y)]^2+24(1-\eta^2)(1-r)}
\le q+\eta H.
\]
The right side is at least one, so squaring is valid. Using
\(q^2=25-24r\), the squared difference is exactly
\[
\begin{aligned}
2\eta\{&
\eta[8(1-y)(4-3x)+12(1-r)]\\
&+(q-1)(8-3x-4y)+8(1-y)\}.
\end{aligned}
\]
Every factor is nonnegative on
\[
0\le\eta\le1,\qquad -1\le x,y,r\le1.
\]
No omitted compatibility condition among \(x,y,r\) can invalidate the
estimate because it holds on the larger box.

For \(u=|\bm b_0\cdot\bm b_1|\), both relevant pairs satisfy \(r\ge-u\).
Moreover,
\[
(5+8u)^2-(25+24u)=8u(7+8u)\ge0,
\]
so
\[
F_0(r)\le\frac7{20}+\frac1{20}\sqrt{25+24u}
\le\frac35+\frac{2u}{5}.
\]

For completeness, the scalar inequalities in the CHSH-deficit estimate have
nonnegative squared differences
\[
\left(1-\frac{u^2}{2}\right)^2-(1-u^2)=\frac{u^4}{4},
\]
\[
\left(2-\frac{u^2}{4}\right)^2-(4-u^2)=\frac{u^4}{16}.
\]
Their right sides are nonnegative on \(0\le u\le1\), so no sign reversal is
introduced by squaring.

\section{The strengthened algebraic strategy}
\label{app:strengthened}

Put
\[
q=\sqrt{7813},\qquad \eta=\frac1q,\qquad
s=\sqrt{1-\eta^2}=\frac{6\sqrt{217}}q.
\]
Use
\[
|\psi_\eta\rangle=
\sqrt{\frac{1+\eta}{2}}|00\rangle+
\sqrt{\frac{1-\eta}{2}}|11\rangle,
\]
Bob observables \(B_0=X,B_1=Z\), and Alice observables
\[
A_0=\frac{6\sqrt{217}\,X+qZ}{125},\qquad
A_1=\frac{6\sqrt{217}\,X-qZ}{125}.
\]
They square to the identity because
\[
36\cdot217+7813=125^2.
\]
The CHSH value is \(S=250/q\).

Set
\[
k=\frac35\sqrt{\frac{q-1}{q+1}}
\]
and
\[
M_0=\frac12\begin{pmatrix}k^2&k\\k&1\end{pmatrix},\qquad
M_1=\frac12\begin{pmatrix}k^2&-k\\-k&1\end{pmatrix},\qquad
M_2=\begin{pmatrix}1-k^2&0\\0&0\end{pmatrix}.
\]
These are positive rank-one effects and sum to the identity. The diagonal
dual operator
\[
\Gamma=
\operatorname{diag}\!\left(
\frac25(1+\eta),\frac6{25}(1-\eta)\right)
\]
dominates all three weighted steered score operators and annihilates the
corresponding effect ranges. Hence
\[
T=\Tr\Gamma=\frac{16+4\eta}{25},
\qquad
10S+T=\frac{16+8q}{25}.
\]

Within this one-parameter family,
\[
f(\eta)=20\sqrt{2-\eta^2}+\frac{16+4\eta}{25}
\]
has
\[
f'(\eta)=-\frac{20\eta}{\sqrt{2-\eta^2}}+\frac4{25},\qquad
f''(\eta)=-\frac{40}{(2-\eta^2)^{3/2}}<0.
\]
Its unique critical point is \(\eta=1/\sqrt{7813}\). This is not a global
POVM optimality claim.

Finally,
\[
(6+80\sqrt{7813})^2-2\cdot5003^2
=-56782+960\sqrt{7813},
\]
and
\[
960^2\cdot7813-56782^2=3\,976\,265\,276>0.
\]
Both sides being compared are positive, these two exact square comparisons
prove \(L_1>U\).

\section{Degenerate cases in the one-binary-party theorem}
\label{app:binary-degeneracies}

If \(B_0,B_1\) are nondegenerate and have linearly independent Bloch axes,
\(\Phi_B\) maps the positive cone onto a three-dimensional Lorentz cone.
An extreme ray of the image lies on the boundary ellipse. Its minimum-norm
Bloch lift in the plane of the two axes has norm one; a perpendicular
component would exceed the Bloch ball. Thus the positive lift in
\(\operatorname{span}_{\R}\{\id,B_0,B_1\}\) is unique and rank one.

If the Bloch axes are parallel or antiparallel, the image is the cone over an
interval and has two extreme rays. Every positive relation reduces to
one-versus-one cancellations or a two-versus-two relation with repeated
zero-padded outcomes. The same purification construction applies.

If one observable is scalar and the other nondegenerate, the image is again
two-dimensional; the scalar input is deterministic and does not add a
coordinate. If both are scalar, Bob is deterministic on both inputs and the
behavior is local. A rank-deficient circuit sum \(\Omega\) forces all lifted
rays to be collinear and is realized by a product purification. Zero effects
receive zero projectors. These cases also show that circuit weights remain
nonnegative and sum to one at the boundary.

\section{Extremal qubit POVMs and the rank-square bound}
\label{app:extremal-povm}

Let \(M_a\) be the nonzero effects of a POVM. For each \(a\), let
\(\mathcal V_a\) be the real vector space of Hermitian operators supported on
\(\operatorname{ran}M_a\). It has dimension \(\rank(M_a)^2\).
Consider the summation map
\[
\bigoplus_a\mathcal V_a\longrightarrow\Herm(2),\qquad
(H_a)_a\longmapsto\sum_aH_a.
\]
If it has a nonzero kernel, then for sufficiently small \(\eps>0\),
\(M_a\pm\eps H_a\ge0\), and the two perturbed tuples are distinct POVMs whose
average is \(M\). Extremality therefore requires injectivity, giving
\[
\sum_a\rank(M_a)^2\le\dim_{\R}\Herm(2)=4.
\]
This necessary condition is all that is used in
\cref{thm:residual-reduction}. Together with linear independence in a
three-dimensional span, it leaves exactly three rank-one effects.

\section{Finite-dimensional POVM duality}
\label{app:duality}

Let
\[
V=\max_{\substack{N_j\ge0\\\sum_jN_j=\id}}
\sum_j\Tr(K_jN_j).
\]
Weak duality says that every \(\Gamma\ge K_j\) obeys
\[
\sum_j\Tr(K_jN_j)\le\Tr\Gamma.
\]
For the reverse direction, consider the closed cone
\[
\mathcal C=
\left\{\left(\sum_jX_j,\sum_j\Tr(K_jX_j)-t\right):
X_j\ge0,\ t\ge0\right\}
\subset\Herm(2)\times\R.
\]
Closedness follows because convergence of \(\sum_jX_j\) bounds every
\(0\le X_j\le\sum_kX_k\).

The point \((\id,V+\eps)\) is outside \(\mathcal C\). Nearest-point
separation provides a nonzero functional
\[
\ell(A,s)=\Tr(\Gamma_\eps A)+\gamma s
\]
nonnegative on \(\mathcal C\) and negative at that point. Since
\((0,-t)\in\mathcal C\), \(\gamma\le0\). If \(\gamma=0\), positivity on all
\((X,\Tr(K_jX))\) would give \(\Gamma_\eps\ge0\), contradicting
\(\Tr\Gamma_\eps<0\). Normalize \(\gamma=-1\). Taking one nonzero \(X_j=X\)
gives \(\Gamma_\eps\ge K_j\), while separation and weak duality give
\[
V\le\Tr\Gamma_\eps<V+\eps.
\]
A bounded convergent subsequence as \(\eps\downarrow0\) yields dual
attainment with \(\Tr\Gamma=V\).

At a primal optimum,
\[
0=\Tr\Gamma-\sum_j\Tr(K_jN_j)
=\sum_j\Tr[(\Gamma-K_j)N_j].
\]
Each term is nonnegative, hence zero. For \(A,B\ge0\),
\(\Tr(AB)=0\) implies \(A^{1/2}BA^{1/2}=0\), and therefore \(AB=0\).
This proves complementary slackness.

\section{Square completion for the incidence Hessian}
\label{app:square-completion}

Let \(x_j=Pr_j\), \(Q=g^{-1}\), and vary
\[
P(t)=P+t\delta P,\qquad g(t)=g+tH
\]
to first order. The inverse expansion is
\[
(g+tH)^{-1}
=Q-tQHQ+t^2QHQHQ+O(t^3).
\]
Substitution in \(F_j=(P r_j)^TQ(P r_j)\) gives
\[
\begin{aligned}
\frac12D^2F_j
={}&(\delta P r_j)^TQ(\delta P r_j)
-2(\delta P r_j)^TQHQ(Pr_j)\\
&+(Pr_j)^TQHQHQ(Pr_j).
\end{aligned}
\]
Writing
\[
\delta P=PW+HQ P
\]
turns the right side into
\[
(PWr_j)^TQ(PWr_j).
\]
Therefore
\[
\sum_j\lambda_jD^2F_j
=2\sum_j\lambda_jr_j^TW^TP^TQPW r_j
=2\Tr(SW\Lambda W^T).
\]
This expansion also fixes the sign convention used at a Bell maximum.

\section{Projective fibers of the quadratic null-ray map}
\label{app:fibers}

Let
\[
\mathcal N_g=\{[x]\in\mathbb{RP}^3:x^Tgx=0\}
\]
for a strict metric \eqref{eq:metric}, and write
\(\Phi(x)=(A,B,C,D)\).

\subsection{Base locus and generic inverse}

If \(x_2=x_3=0\), nullness gives \(x_0x_1=0\), hence \(r_1,r_2\).
If exactly one of \(x_2,x_3\) vanishes, the remaining equations give
\(r_3\) or \(r_4\). If \(x_2x_3\ne0\), all four components of \(\Phi\)
vanish only when
\[
x_0=x_1=-x_2=-x_3,
\]
the ray \(r_5\). Thus the base locus consists exactly of the five coefficient
rays.

The identities
\[
\begin{aligned}
A-B&=x_0(x_2-x_3),&
C-D&=x_1(x_2-x_3),\\
A-C&=x_2(x_0-x_1),&
B-D&=x_3(x_0-x_1),\\
AD-BC&=x_2x_3(x_0-x_1)(x_2-x_3)
\end{aligned}
\]
give
\[
\begin{pmatrix}
(A-B)(A-C)(B-D)\\
(C-D)(A-C)(B-D)\\
(AD-BC)(A-C)\\
(AD-BC)(B-D)
\end{pmatrix}
=\omega(x)x,
\]
where
\[
\omega(x)=x_2x_3(x_2-x_3)(x_0-x_1)^2.
\]
Hence \([x]\) is recovered from \([\Phi(x)]\) whenever \(\omega(x)\ne0\).

\subsection[Exceptional divisor x2=0]{Exceptional divisor \(x_2=0\)}

The target is \([0:B:0:D]\). On the direct branch, scale \(z_3=1\) and
write \((z_0,z_1)=t(B,D)\). Nullness is
\[
t^2BD+2t(bB+dD)=0.
\]
The root \(t=0\) is \(r_4\), leaving at most one nonbase root. On the cross
branch \(z_2\ne0\), target zeros force \(z_0=z_1=-z_3\) and a nonzero target
requires \(B=D\). With \(z_3=1,z_2=t\),
\[
z^Tgz=(1-2(b+d))(1-t).
\]
Strictness \(b+d<1/2\) gives only \(t=1\), the base ray \(-r_5\).

\subsection[Exceptional divisor x3=0]{Exceptional divisor \(x_3=0\)}

The symmetric direct equation is
\[
t^2AC+2t(aA+cC)=0,
\]
and the cross equation is
\[
(1-2(a+c))(1-t)=0.
\]
Again there is at most one nonbase preimage.

\subsection[Exceptional divisor x0=x1]{Exceptional divisor \(x_0=x_1\)}

If \(x_0=x_1=0\), nullness reduces to \(2e x_2x_3=0\). Since \(e>0\),
the point is \(r_3\) or \(r_4\), hence in the base locus. A nonbase point may
therefore be scaled to \(x_0=x_1=1\). Put \(p=x_2,q=x_3\). The null and
target-ratio equations are
\[
N=1+2(a+c)p+2(b+d)q+2epq=0,
\]
\[
R=Bp(1+q)-Aq(1+p)=0.
\]
Their exact resultant in \(p\) is
\[
\operatorname{Res}_p(N,R)
=-(q+1)\left[
B+q\bigl(A(2a+2c-1)+2B(b+d)\bigr)
\right].
\]
At \(q=-1\), strictness gives \(p=-1\), the base ray. The second factor is
not identically zero: if \(B\ne0\) its constant term is nonzero; if \(B=0\),
then \(A\ne0\) and \(2(a+c)-1\ne0\). Once \(q\) is fixed, \(N\) determines
at most one \(p\). The only apparent failure would require
\[
2(a+c)(b+d)=a+b+c+d-\tfrac12,
\]
equivalent to
\[
(2(a+c)-1)(2(b+d)-1)=0,
\]
which strictness excludes.

\subsection[Exceptional divisor x2=x3]{Exceptional divisor \(x_2=x_3\)}

If \(x_2=x_3=0\), nullness gives \(x_0x_1=0\), hence the base rays
\(r_1,r_2\). For a nonbase point scale \(x_2=x_3=1\) and put
\(p=x_0,q=x_1\). The equations are
\[
N=pq+2(a+b)p+2(c+d)q+2e=0,
\]
\[
R=C(p+1)-A(q+1)=0.
\]
The exact resultant is
\[
\operatorname{Res}_p(N,R)
=-(q+1)\left[
Aq+2A(a+b)+2C(c+d)-C
\right].
\]
At \(q=-1\), strictness gives \(p=-1\), the base ray. Strictness also excludes
an identically vanishing second factor or an undetermined null equation.

For a possible cross-branch preimage with \(z_2\ne z_3\), equal target
coordinates give \(z_0=z_1=0\); strict \(e>0\) and nullness then force
\(z_2z_3=0\), again a base ray. These cases exhaust \(\omega=0\), proving
projective injectivity away from the base locus.

\section{Bounded transportation and the rank-zero simulation}
\label{app:transport}

For three ternary labels, let \(c_{ij}=c_{ji}\ge0\), \(c_{ii}=0\), and assume
\[
\sum_{j\ne i}c_{ij}=a_i+b_i,\qquad
\sum_i a_i=\sum_i b_i=\frac12.
\]
Seek \(q_{ij}\) with row sums \(b_i\), column sums \(a_j\), and
\(0\le q_{ij}\le c_{ij}\). Put
\[
q_{ij}=\frac{c_{ij}}2+f_{ij},\qquad f_{ij}=-f_{ji},
\]
\[
d_i=\frac{b_i-a_i}{2},\qquad d_1+d_2+d_3=0.
\]
Set
\[
f_{12}=t,\qquad f_{13}=d_1-t,\qquad f_{23}=d_2+t.
\]
The row and column equations are automatic. Capacity is equivalent to
\[
t\in I_{12}\cap I_{13}\cap I_{23},
\]
where
\[
\begin{aligned}
I_{12}&=[-c_{12}/2,c_{12}/2],\\
I_{13}&=[d_1-c_{13}/2,d_1+c_{13}/2],\\
I_{23}&=[-d_2-c_{23}/2,-d_2+c_{23}/2].
\end{aligned}
\]
The inequalities
\[
|d_i|\le\frac{a_i+b_i}{2}
=\frac{c_{ij}+c_{ik}}2
\]
show that every pair of intervals intersects, including zero-capacity
endpoints. Three closed intervals on the real line with pairwise
intersections have a common point: the largest left endpoint does not exceed
the smallest right endpoint. Any common \(t\) gives the required bounded
transportation matrix.

The deterministic branch construction in
\cref{prop:rank-zero} then reproduces all four input blocks simultaneously.
The supplied rational simulator independently checks the construction on an
interior symmetric instance; generality follows from the interval proof, not
from that one executable example.

\section{Verification boundary}
\label{app:verification}

The exact verifier checks:
\begin{enumerate}[leftmargin=*]
\item dense and sparse Bell coefficients and both explicit strategies;
\item POVM normalization, rank-one factorizations, exact Bell values, and
  discrimination dual certificates;
\item the determinant, discriminant, robust-comparison, deficit, and
  complete-square identities in the projective bound;
\item the Lorentz metric identities, quadratic map, generic inverse, all
  exceptional resultants and factorizations;
\item an exact pure-state conformal-Lorentz instance and the weighted Hessian
  square completion;
\item independence witnesses for the five incidence constraints;
\item bounded-flow identities and an independent rank-zero simulation.
\end{enumerate}

The scripts do not replace the human arguments establishing compactness,
extremal reduction, cone circuits, physical completeness, multiplier signs,
dimension bounds, or exhaustion of cases. Some symbolic checks are exact
regression instances rather than formal proofs of general matrix identities;
the corresponding general identities are derived above.
